% TEMPLATE-MILC-v3.tex - plantilla maestra UNICA de las guias MILC v3.
%
% La usan las tres familias de guias, cada una con su driver:
%   · Grados 6-11  -> scripts/build-guias-g11.py
%   · Bebras       -> scripts/build-guias-bebras.py
%   · Semillero    -> scripts/build-guias-semillero.py
%
% Antes habia tres copias casi identicas (~400 lineas iguales, 6 distintas):
% cada mejora habia que aplicarla tres veces, y en la practica no se hacia
% (el motor de recursos vivia solo en dos de ellas). Lo que varia entre
% familias esta parametrizado en 6 placeholders que cada driver llena
% (se nombran aqui SIN su sintaxis real: el builder reemplaza por texto plano,
% tambien dentro de los comentarios, y un valor multilinea rompe el preambulo):
%
%   HEADER_IZQ        encabezado izquierdo de pagina
%   HEADER_DER        encabezado derecho de pagina
%   PORTADA_ETIQUETA  rotulo sobre el numero grande (GRADO/MOMENTO/MODULO)
%   PORTADA_SERIE     linea de serie de la portada
%   PORTADA_ID        identificador de la guia en la portada
%   PORTADA_CIRCULO   el circulito de grado (°), o vacio si no aplica
%
% El resto de claves las documenta content/guias/_SCHEMA.md. El script valida
% que no queden placeholders sin reemplazar antes de escribir el archivo final.

\documentclass[11pt,letterpaper]{article}
\usepackage[margin=2.0cm]{geometry}
\usepackage[table]{xcolor}
\usepackage{fontspec}
\usepackage[spanish]{babel}
\usepackage{tikz}
% Librerías TikZ para los diagramas de las guías (bloque `recursos:`):
%  · babel  → babel en español vuelve activos caracteres como «>», y TikZ los usa
%             en sus opciones (>=stealth, ->). Sin esta librería, cualquier
%             diagrama con flechas revienta con «\language@active@arg> has an
%             extra }». Debe ir ANTES de que se dibuje nada.
%  · positioning        → sintaxis «below left=2pt and 3pt».
%  · decorations.pathreplacing → llaves ({brace}) para señalar rangos.
\usetikzlibrary{babel,positioning,decorations.pathreplacing}
\usepackage{graphicx}
% Circuitos: el trabajo de electrónica se hace hoy con micro:bit y sus sensores
% integrados (MakeCode, sin cableado), así que no se necesitan esquemáticos.
% Si algún día entran montajes con protoboard, basta descomentar esta línea:
% los diagramas .tex/.tikz declarados en `recursos:` ya se incrustan solos.
% \usepackage{circuitikz}
\usepackage[most]{tcolorbox}
\usepackage{tabularx}
\usepackage{array}
\usepackage{fancyhdr}
\usepackage{titlesec}
\usepackage[document]{ragged2e}  % bandera a la derecha en todo el cuerpo
\usepackage{enumitem}
\usepackage{microtype}

% Helvetica Neue no trae la flecha U+2192, asi que cada «→» escrito literalmente
% en el YAML se compilaba a NADA, en silencio (462 flechas en 61 guias: rutas del
% tipo «paleta → tipografia → lockup» salian rotas). Mapeamos el caracter a su
% version matematica, que si tiene glifo. Asi el autor puede seguir escribiendo
% «→» comodamente en el YAML.
\usepackage{newunicodechar}
\newunicodechar{→}{\ensuremath{\rightarrow}}
% Mismo problema con los simbolos de marca y vinetas: el «✓» de «una palomita ✓»
% salia como cuadrito vacio en el PDF de todas las guias que lo usaban.
\usepackage{pifont}
\newunicodechar{✓}{\ding{51}}
\newunicodechar{✗}{\ding{55}}
\newunicodechar{✔}{\ding{52}}
\newunicodechar{•}{\textbullet}
\newunicodechar{≥}{\ensuremath{\geq}}
\newunicodechar{≤}{\ensuremath{\leq}}

\setmainfont{Helvetica Neue}
\setsansfont{Helvetica Neue}
\setlength{\parindent}{0pt}
\setlength{\parskip}{6pt}
% Sin particion de palabras: en 6.º-7.º una palabra cortada por guion al final
% de linea («Comuni-cacion») frena la lectura mucho mas que una linea corta.
\hyphenpenalty=10000
\exhyphenpenalty=10000
\pretolerance=2000
\tolerance=3000
\setlength{\emergencystretch}{3em}
\linespread{1.10}
\renewcommand{\arraystretch}{1.25}
\setlist{leftmargin=5mm,itemsep=1mm,topsep=1mm}
\newcolumntype{Y}{>{\RaggedRight\arraybackslash}X}

\definecolor{milcMagenta}{HTML}{E6007E}
\definecolor{milcVino}{HTML}{5A0038}
\definecolor{milcTurquesa}{HTML}{0093A5}
\definecolor{milcVerde}{HTML}{58A923}
\definecolor{milcMostaza}{HTML}{E5B400}
\definecolor{milcOcre}{HTML}{B86600}
\definecolor{milcNegro}{HTML}{241816}
\definecolor{milcGris}{HTML}{F7F7F4}
\definecolor{milcLinea}{HTML}{E8E2DD}
\definecolor{milcGradePrimary}{HTML}{0D9488}
\definecolor{milcGradeDark}{HTML}{115E59}
\definecolor{milcGradeSoft}{HTML}{CCFBF1}
\definecolor{milcGradeAccent}{HTML}{CFE7FF}
\definecolor{milcGradeText}{HTML}{FFFFFF}
\color{milcNegro}

\pagestyle{fancy}
\fancyhf{}
\lhead{\small Territorio interior · Educación socioemocional}
\rhead{\small Grado 7° · Momento 2 · MILC v3}
\cfoot{\small\thepage}
\renewcommand{\headrulewidth}{0pt}

\newtcolorbox{titlebox}[1]{
  enhanced,colback=#1,colframe=#1,arc=7mm,boxrule=0pt,
  left=7mm,right=7mm,top=3mm,bottom=3mm,
  before skip=8pt,after skip=8pt,
  fontupper=\sffamily\bfseries\Large\color{white}
}

\newtcolorbox{softbox}[3]{
  enhanced,colback=#2,colframe=#1,borderline west={4pt}{0pt}{#1},
  arc=2mm,boxrule=.4pt,left=4mm,right=4mm,top=3mm,bottom=3mm,
  before skip=6pt,after skip=8pt,title={#3},coltitle=white,
  colbacktitle=#1,fonttitle=\sffamily\bfseries,fontupper=\RaggedRight,
  attach boxed title to top left={xshift=3mm,yshift=-2mm},
  boxed title style={arc=2mm, boxrule=0pt}
}

\newtcolorbox{drawbox}[2]{
  enhanced,colback=white,colframe=#1,arc=4mm,boxrule=1pt,
  left=5mm,right=5mm,top=4mm,bottom=4mm,before skip=8pt,after skip=8pt,
  title={#2},coltitle=white,colbacktitle=#1,fonttitle=\sffamily\bfseries,
  fontupper=\RaggedRight,
  attach boxed title to top left={xshift=4mm,yshift=-2mm},
  boxed title style={arc=3mm, boxrule=0pt}
}

\newtcolorbox{infoband}[1]{
  enhanced,colback=#1!8,colframe=#1!50,arc=2mm,boxrule=.5pt,
  left=4mm,right=4mm,top=2mm,bottom=2mm,before skip=4pt,after skip=4pt,
  fontupper=\small\RaggedRight
}

% Puente narrativo entre secciones (contrato editorial Conéctate).
% Da continuidad: "Acabas de X. Ahora vas a Y porque Z."
% Visualmente: banda izquierda en color del grado, texto en itálica pequeña.
\newtcolorbox{puentebox}{
  enhanced,colback=milcGradeSoft,colframe=milcGradeDark,
  borderline west={3pt}{0pt}{milcGradeDark},
  arc=0mm,boxrule=0pt,left=5mm,right=4mm,top=2.5mm,bottom=2.5mm,
  before skip=4pt,after skip=8pt,
  fontupper=\itshape\small\color{milcNegro}\RaggedRight
}
% La flecha va en modo matematico a proposito: Helvetica Neue NO tiene el glifo
% U+2192, asi que \textrightarrow se compilaba a NADA («Missing character: There
% is no → in font Helvetica Neue») y el puente salia sin flecha. En modo
% matematico la flecha viene de la fuente matematica y si se dibuja.
\newcommand{\puente}[1]{%
  \begin{puentebox}%
    \textcolor{milcGradeDark}{$\rightarrow$}\hspace{2mm}#1%
  \end{puentebox}%
}

\newcommand{\linead}[1]{\noindent\color{milcLinea}\rule{#1}{0.5pt}\color{milcNegro}}
\newcommand{\respuesta}[1]{\par\vspace{2mm}\foreach \n in {1,...,#1}{\noindent\color{milcLinea}\rule{\linewidth}{0.5pt}\par\vspace{5mm}}\color{milcNegro}}
\newcommand{\checkbox}{\textcolor{milcGradeDark}{\fbox{\phantom{x}}}\hspace{5pt}}

% ─── Listas aireadas para las guías socioemocionales ────────────────────
% Componer las verificaciones como «\checkbox texto\\» deja el texto que salta
% de línea debajo del cuadrito y apelmaza el bloque. Estos entornos alinean la
% continuación y airean los ítems. Los usa Territorio Interior; cualquier
% familia puede hacerlo.
\newlist{checklista}{itemize}{1}
\setlist[checklista]{
  label=\textcolor{milcGradeDark}{\fbox{\phantom{x}}},
  leftmargin=9mm, labelsep=3.2mm, itemsep=2.4mm,
  topsep=2.5mm, parsep=0pt, partopsep=0pt, before=\RaggedRight
}
\newlist{pasoslista}{enumerate}{1}
\setlist[pasoslista]{
  label=\textcolor{milcGradeDark}{\sffamily\bfseries\arabic*.},
  leftmargin=8mm, labelsep=3mm, itemsep=2.8mm,
  topsep=1mm, parsep=0pt, partopsep=0pt, before=\RaggedRight
}

\newcommand{\phasecard}[4]{
\begin{tcolorbox}[enhanced,colback=#3,colframe=#2,arc=3mm,boxrule=.5pt,height=3.05cm,valign=center,halign=center,left=1.5mm,right=1.5mm,top=2mm,bottom=2mm]
{\sffamily\bfseries\fontsize{22}{24}\selectfont\textcolor{#2}{#1}}\par
{\sffamily\bfseries\small #4}
\end{tcolorbox}
}

% ── Piezas del rediseno 2026 (legibilidad 6.º-11.º) ───────────────────
% Banda de estacion: reemplaza el titlebox macizo. Titulo a la izquierda,
% dato practico (tiempo/modalidad) a la derecha, en una sola linea.
\newcommand{\milcestacion}[2]{%
\begin{tcolorbox}[enhanced,colback=#1,colframe=#1,arc=2mm,boxrule=0pt,
  left=5mm,right=5mm,top=2.6mm,bottom=2.6mm,before skip=11pt,after skip=7pt]
{\sffamily\bfseries\fontsize{15}{18}\selectfont\color{white}#2}
\end{tcolorbox}}

% Tarjeta de paso numerada: sustituye las tablas «Pilar / Aplicacion», que
% eran formato de consulta docente, no de estudiante. El numero va en un
% circulo sobre el borde para que la secuencia se lea de un vistazo.
\newcommand{\milcpaso}[3]{%
\begin{tcolorbox}[enhanced,colback=white,colframe=milcLinea,arc=1.5mm,boxrule=.9pt,
  left=14mm,right=4mm,top=3mm,bottom=3mm,before skip=4pt,after skip=4pt,
  fontupper=\RaggedRight,
  overlay={\node[anchor=north west,circle,fill=#2,text=white,inner sep=0pt,
    minimum size=7.4mm,font=\sffamily\bfseries\small]
    at ([xshift=3.6mm,yshift=-3.2mm]frame.north west) {#1};}]
#3
\end{tcolorbox}}

% Casilla marcable de tamano util con lapiz (la anterior era \fbox{\phantom{x}},
% de alto variable segun el texto que la seguia).
\newcommand{\milcmarca}{\framebox[3.8mm]{\rule{0pt}{3.4mm}}}

% Fila marcable de lista suelta (checks de la estacion 1).
\newcommand{\milccheckrow}[1]{%
\par\vspace{1.8mm}\noindent
\makebox[6.4mm][l]{\raisebox{-.55mm}{\textcolor{milcNegro}{\milcmarca}}}%
\parbox[t]{\dimexpr\linewidth-6.4mm\relax}{\RaggedRight #1}\par}

% Celda marcable dentro de la rubrica (tabla de autoevaluacion). Misma
% sangria francesa, pero pensada para vivir dentro de una columna X.
\newcommand{\milcopcion}[1]{%
\makebox[5.2mm][l]{\raisebox{-.55mm}{\textcolor{milcNegro}{\milcmarca}}}%
\parbox[t]{\dimexpr\linewidth-5.2mm\relax}{\RaggedRight #1}}

% ── Recursos visuales de la guía ──────────────────────────────────────
% Declarados en el bloque `recursos:` del YAML y emitidos por el builder.
% \guiaFigura   → imagen rasterizada o PDF (foto, carta celeste, montaje).
% \guiaDiagrama → fuente TikZ/circuitikz (.tex) incrustada como vector:
%                 es la vía para circuitos y esquemáticos editables.
% #1 = ruta relativa al .tex · #2 = pie (puede ir vacío)
\newcommand{\guiaFigura}[2]{%
\begin{center}
\includegraphics[width=0.88\linewidth,height=0.38\textheight,keepaspectratio]{#1}\\[1.5mm]
{\footnotesize\itshape\textcolor{milcNegro!75}{#2}}
\end{center}
}

\newcommand{\guiaDiagrama}[2]{%
\begin{center}
\input{#1}\\[1.5mm]
{\footnotesize\itshape\textcolor{milcNegro!75}{#2}}
\end{center}
}

\begin{document}
\thispagestyle{empty}

% ── Portada ───────────────────────────────────────────────────────────
% Antes: un rectangulo de color a pagina completa. Se veia bien en pantalla,
% pero en la fotocopiadora del colegio se comia el toner y salia gris sucio.
% Ahora la banda ocupa el tercio superior y el resto va en blanco.
\begin{tikzpicture}[remember picture,overlay]
  \path[fill=milcGradePrimary]
    (current page.north west) rectangle ([yshift=-10.6cm]current page.north east);

  \node[anchor=north west,text=milcGradeText] at ([xshift=2.0cm,yshift=-1.55cm]current page.north west)
    {\sffamily\bfseries\fontsize{12}{14}\selectfont MOMENTO};
  \node[anchor=north west,text=milcGradeText,opacity=.85] at ([xshift=2.0cm,yshift=-2.35cm]current page.north west)
    {\sffamily\bfseries\fontsize{8.5}{10}\selectfont TERRITORIO INTERIOR · SOCIOEMOCIONAL};
  \node[anchor=north east,text=milcGradeText,opacity=.85,align=right] at ([xshift=-2.0cm,yshift=-1.55cm]current page.north east)
    {\sffamily\bfseries\fontsize{9}{12}\selectfont I.E. Sor María Juliana\\Cartago, Valle del Cauca};

  \node[anchor=west,text=milcGradeText] (gradonum) at ([xshift=1.85cm,yshift=-7.1cm]current page.north west)
    {\sffamily\bfseries\fontsize{112}{112}\selectfont 2};
  % El circulito de grado (°) solo aplica a las guias de grado («11°»); en
  % Bebras y Semillero el numero grande es un momento/modulo, no un grado.
  % Se posiciona relativo al nodo `gradonum`, no en coordenadas absolutas.
  

  \node[anchor=west,text=milcGradeText,text width=11.4cm,align=left] at ([xshift=1.5cm,yshift=0.5cm]gradonum.east)
    {
     {\sffamily\bfseries\fontsize{9}{11}\selectfont\MakeUppercase{Territorio interior · Socioemocional}}\\[3mm]
     {\sffamily\bfseries\fontsize{21}{25}\selectfont\setlength{\baselineskip}{27pt}La amistad\\[4pt]no es propiedad\\[4pt]querer no es tener}\\[3.5mm]
     {\sffamily\bfseries\fontsize{15}{18}\selectfont Grado 7° · Momento 2}
    };
\end{tikzpicture}

\vspace*{9.4cm}
\vfill

{\sffamily\bfseries\fontsize{8}{10}\selectfont\textcolor{milcGradeDark}{LO QUE TE LLEVAS HOY}}

\begin{tcolorbox}[enhanced,colback=white,colframe=milcNegro,arc=2mm,boxrule=1.6pt,
  left=5mm,right=5mm,top=4mm,bottom=4mm,before skip=3pt,after skip=12pt,
  fontupper=\RaggedRight\fontsize{11}{15.5}\selectfont]
Tu mapa del círculo: quiénes están hoy en tu vida y con qué lazo, dónde aparecen los celos y qué hacen contigo, y un acto concreto de apertura —invitar a alguien, soltar una exigencia— con fecha.
\end{tcolorbox}

\begin{tcolorbox}[enhanced,colback=milcGris,colframe=milcGris,arc=2mm,boxrule=0pt,
  left=5mm,right=5mm,top=3.5mm,bottom=3.5mm,after skip=12pt]
{\sffamily\bfseries\fontsize{8}{10}\selectfont\textcolor{milcNegro!55}{ESTA GUÍA ES DE}}\\[3mm]
\begin{tabularx}{\linewidth}{X p{3.2cm} p{3.2cm}}
\linead{\linewidth} & \linead{3.0cm} & \linead{3.0cm}\\[-1mm]
{\fontsize{8.5}{10}\selectfont\textcolor{milcNegro!55}{Nombre completo}} &
{\fontsize{8.5}{10}\selectfont\textcolor{milcNegro!55}{Grupo}} &
{\fontsize{8.5}{10}\selectfont\textcolor{milcNegro!55}{Fecha}}\\
\end{tabularx}
\end{tcolorbox}

\vfill

\noindent\textcolor{milcLinea}{\rule{\linewidth}{1pt}}\\[2mm]
{\sffamily\bfseries\fontsize{9.5}{12}\selectfont Autor: Dr. Álvaro Cárdenas Orozco}\\[1mm]
{\sffamily\fontsize{8.5}{11}\selectfont\textcolor{milcNegro!55}{Metodología MILC · Apertura ancestral · Vínculos y celos · Triángulo Dussel-Estoicismo-Floridi}}

\newpage

\section*{Apertura: La amistad no es propiedad: querer a alguien no es tenerlo}

\begin{softbox}{milcGradeDark}{milcGradeSoft}{Saber ancestral}
En buena parte de Colombia hay un parentesco que no viene con la sangre: \textbf{se elige}. Es el \textbf{compadrazgo}. Nace del bautizo: unos padres escogen a alguien para que apadrine a su hijo, y desde ese día esas dos familias quedan unidas ---los padres y los padrinos se llaman \textbf{compadres}---, con obligaciones reales de apoyo mutuo que se extienden incluso a los hijos del padrino, que quedan como \emph{hermanos espirituales}. \textbf{Y ahora fíjate en las dos reglas que tiene, porque enseñan todo:} la primera es inflexible ---\textbf{unos padres no pueden ser padrinos de su propio hijo}---; la segunda es una costumbre muy vieja ---\textbf{los padrinos no suelen ser parientes}---. \emph{Es decir: el lazo está diseñado a propósito para traer gente de afuera.} Ninguna familia se apadrina a sí misma. \textbf{Lo que se busca no es cerrar el círculo, sino ampliarlo:} que el hijo tenga más adultos que respondan por él, no menos. \emph{Piensa en eso la próxima vez que sientas que un amigo tuyo «te está reemplazando» porque hizo otro amigo.}
\end{softbox}

\begin{softbox}{milcTurquesa}{milcGris}{Saber contemporáneo}
Los \textbf{celos de amistad} son de las cosas más comunes y menos habladas de la secundaria: la sensación de que tu mejor amiga se está yendo con otra, de que ya no te buscan, de que sobras. \textbf{Sentirlo es normal.} Lo que la investigación agrega es incómodo y vale la pena saberlo: cuando un equipo dirigido por Jeffrey Parker midió los celos de amistad en casi 500 adolescentes, encontró que \textbf{quienes tenían más fama de celosos también se comportaban de forma más agresiva} y tenían más dificultades con el grupo ---y que tanto los celos que uno mismo reconoce como los que le ven los demás \textbf{terminan asociados a la soledad}--- (Parker y otros, 2005). \emph{Léelo despacio, porque es una trampa perfecta:} los celos existen para no perder a alguien, y lo que suelen producir es exactamente lo que temen. \textbf{La distinción que salva es sencilla:} los \textbf{celos} son una emoción, llegan solos y no se eligen; el \textbf{control} ---exigir, reclamar, revisar, castigar con silencio--- es una conducta, y esa sí se elige.
\end{softbox}

\begin{softbox}{milcMostaza}{milcGris}{Pregunta-puente}
\textit{El compadrazgo existe para que un niño tenga \textbf{más gente} que responda por él, no para que tenga una sola. \textbf{¿En qué momento empezaste a creer que querer a alguien significaba ser el único}\ldots y qué te ha costado esa idea?}
\end{softbox}

\begin{softbox}{milcVerde}{milcGris}{Saber y saber hacer}
Al terminar podrás: (1) \textbf{identificar} las situaciones concretas en que aparecen tus celos y qué hiciste con ellos; (2) \textbf{explicar} la diferencia entre celos y control, y qué produce cada uno según la evidencia; (3) \textbf{analizar} tu propio círculo: quién está, con qué lazo y a quién has dejado por fuera; (4) \textbf{crear} un acto concreto de apertura y traducir una exigencia tuya en una petición.
\end{softbox}



\small
\begin{tabularx}{\textwidth}{>{\bfseries}p{3.0cm}Y}
\rowcolor{milcGradePrimary!12}Autor & Dr. Álvaro Cárdenas Orozco\\
DBA & Comprende los fenómenos que afectan a su territorio, regula sus emociones ante ellos y acompaña a otros con empatía (transversal 6°--11°, articulado con la política «Escuela, territorio de vida» del MEN, Resolución 006519 de 2025, y la Ley 1620 de 2013)\\
Referentes & Primera ayuda psicológica (OMS, 2011/2012) · Directrices IASC (2007) · USGS y Servicio Geológico Colombiano · Pueblo nasa y la minga (CRIC) · Dussel · Estoicismo · Floridi\\
Producto final & Tu mapa del círculo: quiénes están hoy en tu vida y con qué lazo, dónde aparecen los celos y qué hacen contigo, y un acto concreto de apertura —invitar a alguien, soltar una exigencia— con fecha.\\
\end{tabularx}
\normalsize

\puente{En el momento anterior trabajaste el escuchar. Hoy vas a lo que aparece cuando el otro también escucha a alguien más: \textbf{los celos}, y lo que uno hace con ellos.}

\milcestacion{milcGradeDark}{Ruta MILC: así vamos a aprender}
\begin{tabularx}{\textwidth}{>{\centering\arraybackslash}X>{\centering\arraybackslash}X>{\centering\arraybackslash}X>{\centering\arraybackslash}X}
\phasecard{1}{milcVerde}{milcGris}{Escucha\\[-1mm]\small Observo, escucho y pregunto.} &
\phasecard{2}{milcTurquesa}{milcGris}{Sistematización\\[-1mm]\small Distingo y reviso.} &
\phasecard{3}{milcMagenta}{milcGris}{Praxis\\[-1mm]\small Construyo y aplico.} &
\phasecard{4}{milcVino}{milcGris}{Evaluación\\[-1mm]\small Reflexión liberadora.}
\end{tabularx}


\puente{Primero tus casos, en privado. No se trata de acusarte de celoso: se trata de ver dónde aparece y qué haces cuando aparece.}

\milcestacion{milcVerde}{Estación 1 de 3 · Escucha}

\begin{softbox}{milcVerde}{milcGris}{Escena para escuchar}
\textbf{Actividad 1 · IDENTIFICA --- Dónde aparecen mis celos.} \textbf{Nadie va a leer esta página.} Escribe \textbf{tres situaciones} de este año en que sentiste que alguien que te importa se estaba yendo con otra persona: en el salón, en el barrio, en la casa ---también hay celos con hermanos---, en un grupo de chat. Para cada una, tres columnas: \textbf{qué pasó exactamente} (los hechos, no la interpretación), \textbf{qué sentiste} (con la palabra más precisa que encuentres: celos, miedo, tristeza, rabia, vergüenza) y \textbf{qué hiciste}. \textbf{Sé honesto con la tercera columna:} dejar de hablarle, contestar cortante, hacer un comentario, hablar de esa persona con otros, mirar sus historias veinte veces, «no hacer nada» pero quedarse rumiando. \emph{Y una línea final: ¿qué pasó después con esa amistad? ¿Se acercó, se alejó, quedó igual?}
\end{softbox}

{\sffamily\bfseries\fontsize{8}{10}\selectfont\textcolor{milcNegro!55}{MARCA CUANDO LO HAYAS HECHO}}
\milccheckrow{Escribiste tres situaciones reales, separando los hechos de tu interpretación.}
\milccheckrow{Nombraste lo que sentiste con una palabra precisa, y no solo «celos».}
\milccheckrow{Escribiste con honestidad qué hiciste, incluido lo que hiciste en silencio.}

\begin{infoband}{milcVerde}


\textbf{En tu cuaderno (Actividad 1 · IDENTIFICA).} Título: «Actividad 1. Dónde aparecen mis celos». \textbf{Formato:} tres filas con tres columnas (qué pasó · qué sentí · qué hice) + una línea sobre qué pasó después con esa amistad. \textbf{Extensión:} media página. \textbf{Sabes que terminaste cuando} en la primera columna hay \textbf{hechos} ---«se sentó con ella en el descanso»--- y no conclusiones ---«ya no le importo»---. Esta página es tuya: \textbf{nadie más tiene que leerla si tú no quieres}.
\end{infoband}

\puente{Ya los tienes ubicados. Ahora la distinción que lo cambia todo, y el dato incómodo sobre a dónde lleva el control.}

\milcestacion{milcTurquesa}{Fase 2 · Sistematización: entender los celos}

\begin{softbox}{milcTurquesa}{milcGris}{Cuatro claves sobre querer sin poseer}
Cuatro claves para que los celos dejen de manejar tus amistades sin que tú te des cuenta.
\end{softbox}

\milcpaso{1}{milcTurquesa}{\textbf{Celos y control son dos cosas, y solo la segunda te define.} Los \textbf{celos} son una emoción: llegan solos, avisan que algo te importa y no se eligen ---igual que la rabia o el miedo---. \textbf{No hay que sentirse mala persona por sentirlos.} El \textbf{control} es lo que uno hace con ellos: exigir que no hable con alguien, revisar con quién anda, castigar con silencio, hacer sentir culpable, poner a elegir. \emph{Nadie puede pedirte cuentas por lo que sientes; sí por lo que haces.} La frase que separa las dos: \emph{«siento celos» está bien; «no te juntes con ella» es otra cosa}.}
\milcpaso{2}{milcTurquesa}{\textbf{Y esto es lo que producen, medido.} Parker y su equipo estudiaron los celos de amistad en casi 500 adolescentes: los que tenían más fama de celosos \textbf{se comportaban de forma más agresiva}, tenían más problemas con el grupo, y ---aquí está lo duro--- los celos aparecían asociados con \textbf{más soledad} (Parker y otros, 2005). \emph{Es la trampa completa:} los celos existen para no perder a alguien, y el control que producen es justamente lo que aleja a la gente. \textbf{Nadie se queda donde lo vigilan.}}
\milcpaso{3}{milcTurquesa}{\textbf{Un vínculo sano abre el círculo; uno celoso lo cierra.} Vuelve al compadrazgo: los padres \emph{no} pueden apadrinar a su propio hijo, y los padrinos \emph{no} suelen ser parientes. La tradición eligió, a propósito, \textbf{traer a alguien de afuera}. Porque la pregunta que importa no es «¿quién es más importante para él?», sino \textbf{«¿cuánta gente lo está cuidando?»}. \emph{Cuando tu amiga hace otra amiga, no perdiste la mitad de tu amiga: ella ganó a alguien más, y eso, si de verdad te importa, es una buena noticia.}}
\milcpaso{4}{milcTurquesa}{\textbf{Lealtad no es exclusividad ---y la diferencia se nota en los hechos---.} Un amigo leal te defiende cuando no estás, te dice la verdad, cumple lo que promete y no cuenta lo que le contaste. \textbf{Nada de eso requiere que sea tu único amigo.} La exclusividad, en cambio, no mide nada: alguien puede estar solo contigo y tratarte mal. \emph{Confundir las dos hace que le exijas a la gente la prueba equivocada:} pides que no tenga a nadie más, cuando lo que en realidad necesitas es saber que puedes confiar.}

\begin{drawbox}{milcTurquesa}{De la exigencia a la petición: cómo se traduce}
\begin{pasoslista}
\item \textbf{La exigencia} habla del otro y le quita opciones: «no te sientes con ella», «tienes que contestarme ya», «o ella o yo».
\item \textbf{La petición} habla de ti y deja libre al otro: «me sentí por fuera cuando se fueron sin avisarme», «me gustaría que un día de la semana nos veamos solos».
\item \textbf{La prueba para distinguirlas:} si tu frase empieza con «tú siempre» o «tú nunca», es exigencia. Si empieza con «yo sentí» o «me gustaría», es petición.
\item \textbf{Y lo que casi nadie hace:} decir la petición \emph{en vez} de castigar con silencio. El silencio no comunica nada; solo alarga.
\end{pasoslista}
\end{drawbox}

\begin{softbox}{milcMostaza}{milcGris}{Errores comunes y su corrección}
\textbf{«Si le importara, no haría eso»:} interpretar no es saber; casi siempre falta información. \textbf{Castigar con silencio:} el otro no adivina la razón, solo siente el frío y se aleja. \textbf{Pedir pruebas de lealtad:} el que necesita pruebas todo el tiempo no está construyendo confianza, la está gastando. \textbf{Hablar del amigo con terceros en vez de con él:} eso no resuelve y suma gente al conflicto. \textbf{Confundir celos con intuición:} a veces uno acierta, pero eso se comprueba preguntando, no vigilando. \textbf{Creerse mala persona por sentir celos:} la emoción no es el problema; lo que se hace con ella, sí.
\end{softbox}

\begin{infoband}{milcTurquesa}


\textbf{En tu cuaderno (Actividad 2 · EXPLICA).} Título: «Actividad 2. Celos, control y lealtad». \textbf{Formato:} la diferencia entre celos y control con un ejemplo tuyo; qué encontró el estudio de Parker en dos líneas; y \textbf{una exigencia tuya traducida a petición}, con las dos versiones escritas. \textbf{Extensión:} media página. \textbf{Sabes que terminaste cuando} tu petición no contiene ningún «tú siempre» ni ningún «tú nunca».
\end{infoband}

\puente{Ya lo entiendes. Es hora de dibujar tu círculo y hacer algo con él ---abrirlo, que es más difícil que decirlo---.}

\milcestacion{milcMagenta}{Fase 3 · Praxis: abrir el círculo}

\begin{softbox}{milcMagenta}{milcGris}{Cómo se dibuja y se abre un círculo}
Ahora se dibuja. Un círculo no es una lista de mejores amigos: es el mapa de la gente que hay en tu vida y del lazo que tienes con cada uno ---y de los huecos---.
\end{softbox}

\milcpaso{1}{milcMagenta}{\textbf{Dibujar el mapa (12 min).} Pon un punto en el centro: tú. Alrededor, \textbf{tres anillos}: el de todos los días, el de cada tanto, y el de los que están aunque no se vean. Ubica ahí a tu gente ---amigos, familia, vecinos, primos, alguien del barrio, un adulto de confianza--- y \textbf{anota el lazo}: sangre, elección, cercanía, historia. \emph{Fíjate en dos cosas: cuánta gente hay en total, y cuántos de tus lazos son de los que se eligen ---como el compadrazgo--- y no de los que vienen dados.}}
\milcpaso{2}{milcMagenta}{\textbf{Marcar los celos (8 min).} Con otro color, marca \textbf{dónde} sientes celos en ese mapa: con quién, y \emph{de quién}. A veces sorprende: no siempre es del amigo, a veces es del que llegó nuevo, o de un hermano. \emph{Marcar no es acusarte: es ubicar el lugar exacto donde tu círculo se está apretando.}}
\milcpaso{3}{milcMagenta}{\textbf{Traducir la exigencia (10 min).} Toma la exigencia que escribiste en la Actividad 2 y \textbf{decide si la vas a decir}, y a quién. Escríbela ya lista, como se la dirías: corta, en primera persona, sin «siempre» ni «nunca». \textbf{Ensáyala con tu pareja de trabajo}, en voz alta ---suena distinto dicha que escrita, y hay que oírla antes---. \emph{Si al decirla en voz alta te suena a reclamo, todavía no está traducida.}}
\milcpaso{4}{milcMagenta}{\textbf{El acto de apertura (10 min).} Elige \textbf{una} acción concreta, con fecha, que \textbf{amplíe} tu círculo en vez de cerrarlo: invitar al que siempre queda por fuera, sentarte con alguien nuevo, presentarle tu amiga a otra amiga sin sentirte en peligro, escribirle a alguien con quien te distanciaste por celos. \emph{Que le sirva a alguien más, no solo a ti para sentirte mejor.} \textbf{Escríbela con día.}}

\begin{drawbox}{milcMagenta}{Antes de cerrar tu mapa}
\begin{checklista}
\item Tu mapa tiene los tres anillos y anota el \textbf{lazo} de cada persona, no solo el nombre.
\item Marcaste con otro color dónde aparecen los celos, y de quién.
\item Tienes una exigencia traducida a petición, escrita como se la dirías y \textbf{ensayada en voz alta}.
\item Tu acto de apertura tiene \textbf{día} y le sirve a alguien más, no solo a ti.
\item Identificaste al menos una persona que has dejado por fuera de tu círculo sin darte cuenta.
\item Sabes decir en qué se nota la lealtad de alguien, sin nombrar la exclusividad.
\end{checklista}
\end{drawbox}

\begin{softbox}{milcMagenta}{milcGris}{Plantilla de trabajo}
\textbf{Plantilla del mapa.} \emph{Centro:} yo · \emph{Anillo 1 (todos los días):} \_\_\_\_ · \emph{Anillo 2 (cada tanto):} \_\_\_\_ · \emph{Anillo 3 (están aunque no se vean):} \_\_\_\_ · \emph{Lazo de cada uno:} sangre / elección / cercanía / historia. \\[1mm]
\textbf{Traducción.} \emph{Exigencia:} «\_\_\_\_». \emph{Petición:} «Me sentí \_\_\_\_ cuando \_\_\_\_. Me gustaría \_\_\_\_». \\[1mm]
\textbf{Acto de apertura.} \emph{«Voy a \_\_\_\_, el día \_\_\_\_, y le sirve a \_\_\_\_».}
\end{softbox}

\begin{infoband}{milcMagenta}


\textbf{En tu cuaderno (Actividad 3 · CREA).} Título: «Actividad 3. Mi círculo». \textbf{Formato:} el mapa dibujado con sus tres anillos y los lazos + los celos marcados + la petición escrita + el acto de apertura con día. \textbf{Extensión:} una página. \textbf{Sabes que terminaste cuando} tu mapa muestra \textbf{a alguien que podrías invitar} y no solo a los que ya están.
\end{infoband}

\puente{El producto es ese mapa y un acto con fecha. \emph{«Voy a ser menos celoso» no es un acto: es una intención sin cuerpo.}}

\milcestacion{milcMostaza}{Producto de la sesión}
\textbf{Mapa del círculo: anillos, lazos, celos ubicados, una petición traducida y un acto de apertura con fecha}.
\begin{softbox}{milcMostaza}{milcGris}{Criterios de calidad}
(1) El mapa tiene tres anillos y registra el tipo de lazo de cada persona. (2) Los celos están ubicados en el mapa: con quién y de quién. (3) Hay una exigencia traducida a petición, en primera persona y sin «siempre» ni «nunca». (4) El acto de apertura es concreto, tiene día y beneficia a otra persona. (5) Se distingue por escrito lealtad de exclusividad, con señales observables de la primera.
\end{softbox}

\puente{Antes de cerrar, revisa si tu acto de apertura le sirve a alguien más ---no solo a ti para sentirte mejor---.}

\milcestacion{milcVino}{Evaluación liberadora · cinco dimensiones}

\begin{softbox}{milcVino}{milcGris}{Sentido de la evaluación}
Evaluar no es solo revisar una nota. Es reconocer qué aprendí, qué cambió en mí, qué emociones manejé, cómo me relaciono con la ciudad y la comunidad, y qué le dejaré a quien viene detrás.
\end{softbox}

\textbf{Mi reflexión liberadora en cinco dimensiones.} Completa cada frase iniciada:

\small
\begin{tabularx}{\textwidth}{>{\bfseries\RaggedRight\arraybackslash}p{3.4cm}Y>{\RaggedRight\arraybackslash}p{4.6cm}}
\rowcolor{milcVino}\color{white}Dimensión & \color{white}\bfseries Mi respuesta (frame starter) & \color{white}\bfseries Algo que descubrí\\
1. Personal & Lo que más cambió en mí fue \linead{6.5cm} & \linead{4.2cm}\\
2. Emocional & La emoción más fuerte fue \linead{2.5cm} y la regulé \linead{3cm} & \linead{4.2cm}\\
3. Ciudadana & En democracia esto importa porque \linead{6cm} & \linead{4.2cm}\\
4. Local (Cartago) & En mi barrio sería útil para \linead{6.5cm} & \linead{4.2cm}\\
5. Intergeneracional & A mi abuela/abuelo le diría \linead{6.5cm} & \linead{4.2cm}\\
\end{tabularx}
\normalsize

\begin{infoband}{milcVino}
\textbf{Para la próxima sustentación.} Vuelve a esta tabla en seis meses. Verás que las respuestas cambian — no porque lo aprendido cambie, sino porque tú cambias. La evaluación liberadora es una conversación contigo a través del tiempo.
\end{infoband}


\puente{Cerramos con tres voces: lo que hay más allá del círculo cerrado, lo que sí depende de ti en un vínculo, y qué le hace a un curso que los círculos se cierren.}

\milcestacion{milcGradeDark}{Triángulo de pensamiento · cierre filosófico}

\begin{softbox}{milcVino}{milcGris}{Dussel · lente del nosotros}
\textit{«Analéctico quiere indicar el hecho real humano por el que todo hombre, todo grupo o pueblo se sitúa siempre más allá (aná-) del horizonte de la totalidad.»}

{\footnotesize\itshape\color{milcNegro!70}--- Enrique Dussel · Filosofía de la liberación (1977), §2.4}

Un círculo cerrado ---«nosotros dos y nadie más»--- es lo que Dussel llamaría una \textbf{totalidad}: un mundo que se cree completo y que trata al que llega como una amenaza. \emph{Lo analéctico es lo contrario: reconocer que siempre hay alguien más allá del borde}, y que ese alguien no viene a quitar nada. \textbf{El compadrazgo es analéctico en estado puro:} obliga a salir de la familia para buscar afuera a quien va a cuidar del hijo. \emph{Cada vez que un grupo del salón se cierra, alguien queda «a la vera del camino»; y casi siempre es el mismo.}

\textbf{¿A quién estoy tratando como amenaza cuando en realidad es alguien más que llegó a cuidar lo mismo que yo cuido?}
\end{softbox}
\linead{\linewidth}\\[1mm]
\linead{\linewidth}

\begin{softbox}{milcMostaza}{milcGris}{Estoicismo · Epicteto · Enquiridión, 1 (c. 125 d.C.; trad. desde E. Carter) · cuidado interior}
\textit{«Algunas cosas están en nuestro poder y otras no. Están en nuestro poder la opinión, el impulso, el deseo, el rechazo; en una palabra, todo aquello que es obra nuestra.»}

{\footnotesize\itshape\color{milcNegro!70}--- Epicteto · Enquiridión, 1 (c. 125 d.C.; trad. desde E. Carter)}

A quién quiere tu amiga, con quién se sienta, a quién le escribe: \textbf{nada de eso está en tu poder}, y todo el esfuerzo que pongas en controlarlo es esfuerzo perdido ---y, según los datos, esfuerzo que la aleja---. \textbf{Sí está en tu poder} cómo la tratas, si dices lo que sentiste en vez de castigar con silencio, y si eres la clase de amigo por el que alguien quiere volver. \emph{Lo que retiene a la gente no es el control: es que estar contigo sea bueno.}

\textbf{De todo lo que hago cuando siento celos, ¿qué parte busca controlar al otro y qué parte me hace mejor amigo?}
\end{softbox}
\linead{\linewidth}\\[1mm]
\linead{\linewidth}

\begin{softbox}{milcTurquesa}{milcGris}{Floridi · lente de la infoesfera}
\textit{«El florecimiento de las entidades informacionales y de toda la infoesfera debe promoverse preservando, cultivando y enriqueciendo sus propiedades.»}

{\footnotesize\itshape\color{milcNegro!70}--- Luciano Floridi · La ética de la información (2013; trad.)}

Un salón también florece o se empobrece, y lo deciden cosas pequeñas: quién se sienta con quién, a quién invitan al grupo del chat, quién almuerza solo. \textbf{Un curso hecho de círculos cerrados es un curso pobre}, aunque cada círculo por dentro se sienta cálido: siempre hay alguien afuera, y ese alguien aprende peor, viene menos y se enferma más. \emph{Enriquecer, en tu curso, es tan concreto como correrse un puesto para que quepa alguien más.}

\textbf{En mi salón, ¿quién come solo? ¿Y qué me costaría a mí, en concreto, que no comiera solo mañana?}
\end{softbox}
\linead{\linewidth}\\[1mm]
\linead{\linewidth}

\begin{infoband}{milcNegro}
\textbf{Nota para el docente.} Las tres son citas textuales y llevan su referencia debajo. Puedes remitir a la obra en clase.
\end{infoband}

\milcestacion{milcVino}{Mi autoevaluación · marca una casilla por fila}

{\small
\setlength{\extrarowheight}{2.2pt}
\begin{tabularx}{\textwidth}{>{\RaggedRight\arraybackslash\bfseries}p{3.0cm}YYY}
\rowcolor{milcVino}\color{white}\bfseries Criterio & \color{white}\bfseries Lo logré & \color{white}\bfseries En proceso & \color{white}\bfseries Necesito apoyo\\[.6mm]
Comprensión del tema & \milcopcion{Explico las ideas con mis palabras.} & \milcopcion{Reconozco las ideas, falta precisión.} & \milcopcion{Necesito repasar lo básico.}\\[.8mm]
\rowcolor{milcGris}Distinguir celos de control & \milcopcion{Separo lo que siento de lo que exijo, y reconozco qué produce cada uno.} & \milcopcion{Reconozco los celos, pero todavía justifico la exigencia.} & \milcopcion{Creo que si un amigo tiene otros amigos, me está fallando.}\\[.8mm]
Aplicación y producto & \milcopcion{Mi producto cumple los criterios.} & \milcopcion{Está armado, con ajustes pendientes.} & \milcopcion{Aún está en construcción.}\\[.8mm]
\rowcolor{milcGris}Reflexión 5 dimensiones & \milcopcion{Respondí cada una con honestidad.} & \milcopcion{Respondí algunas con detalle.} & \milcopcion{Mis respuestas son muy generales.}\\[.8mm]
Triángulo de pensamiento & \milcopcion{Conecté las 3 voces con mi trabajo.} & \milcopcion{Conecté al menos una.} & \milcopcion{No las relaciono con mi proceso.}\\[.8mm]
\end{tabularx}}

\vspace{3mm}
\textbf{Hoy entendí que:}\\[1mm]
\linead{\linewidth}\\[1mm]
\linead{\linewidth}

\textbf{Mi compromiso para la próxima sesión es:}\\[1mm]
\linead{\linewidth}\\[1mm]
\linead{\linewidth}

\begin{softbox}{milcMostaza}{milcGris}{Cita de despedida · Marco Aurelio}
\textit{``El obstáculo en el camino se vuelve el camino. Lo que estorba al camino se vuelve camino.''}\\[1mm]
Cada error de hoy, bien observado, se convierte en pieza del camino que pisarás mañana.
\end{softbox}

\small
\begin{tabularx}{\textwidth}{>{\bfseries}p{3.6cm}Y}
\rowcolor{milcGradeDark!12}Nombre completo: & \linead{10cm}\\
Fecha: & \linead{4cm} \quad Curso/grupo: \linead{4cm}\\
Mi firma: & \linead{9cm}\\
\end{tabularx}
\normalsize

\begin{infoband}{milcGradeDark}
\textbf{Para la próxima sesión.} Dos cosas. \textbf{Primero, haz tu acto de apertura} el día que escribiste, y anota qué pasó ---incluido si te dio pena, si salió raro o si la otra persona reaccionó distinto a lo que esperabas---. \emph{Y si te atreviste a decir la petición, anota también cómo se sintió decirla en vez de guardártela.} \textbf{Segundo, pregunta por los padrinos de tu casa:} averigua quiénes son o fueron tus padrinos o los de alguien de tu familia, \textbf{cómo los escogieron} y si esa relación se mantuvo con los años. Pregunta específicamente si eran parientes o no, y por qué eligieron a esa persona. \textbf{Trae:} qué pasó con tu acto de apertura y la historia de un compadrazgo de tu familia. \emph{Vas a encontrar que casi siempre se escogió a alguien por una razón muy concreta: que fuera responsable, que viviera cerca, que quisiera de verdad al niño. Nadie eligió padrino por exclusividad.}
\end{infoband}

\end{document}
